\subsubsection{Particula En tubo}

Una partícula se mueve por el interior de un tubo de largo $2R$ que gira con una velocidad angular constante $\omega_0$. La partícula inicia su movimiento desde el punto medio del tubo, desplazándose por su interior con una rapidez constante $v_0$ respecto al mismo. Determine:

\begin{enumerate}
    \item El radio de curvatura de la trayectoria descrita, en función del tiempo.
    \item La distancia recorrida por la partícula desde que inicia su movimiento hasta que llega al extremo del tubo.
\end{enumerate}

\begin{figure}[htbp]
    \centering
    \begin{tikzpicture}[>=Latex, scale=1.5]
        
        % --- Parámetros ---
        \def\Ang{35}       % Ángulo de inclinación del tubo
        \def\TubeL{4.5}    % Longitud del tubo (2R)
        \def\TubeW{0.5}    % Ancho transversal del tubo
        \def\PartPos{2.5}  % Posición radial de la partícula

        % --- 1. Ejes de Referencia ---
        \draw[->, thick, gray!80] (-0.8,0) -- (5.5,0) node[right, text=black] {$x$};
        \draw[->, thick, gray!80] (0,-0.8) -- (0,4) node[above, text=black] {$y$};

        % --- 2. Tubo con Sombreado y Borde Redondeado ---
        \begin{scope}[rotate=\Ang]
            % Trazado del tubo: comienza con un arco para redondear el origen
            \shadedraw[top color=gray!50, bottom color=gray!50, middle color=white, draw=black, thick] 
                (0, \TubeW/2) arc (90:270:\TubeW/2) -- (\TubeL, -\TubeW/2) -- (\TubeL, \TubeW/2) -- cycle;
            
            % Línea de centro constructiva
            \draw[dash dot, gray!60, thin] (0,0) -- (\TubeL,0);

            % --- 3. Partícula y Vectores Relativos ---
            % Esfera de la partícula
            \shade[ball color=gray!30, draw=black, thin] (\PartPos, 0) circle (0.15);
            \node[below right, yshift=-3pt, font=\small] at (\PartPos, 0) {$m$};
            
            % Vector velocidad v0 (Azul)
            \draw[->, thick, blue] (\PartPos+0.18, 0) -- (\PartPos+1.3, 0) node[below, font=\small] {$v_0$};
            
            % Cota del largo 2R
            \draw[<->, thin] (0, \TubeW/2 + 0.25) -- (\TubeL, \TubeW/2 + 0.25) node[midway, above, font=\small] {$2R$};
        \end{scope}

        % --- 4. Articulación / Pivote Mejorado ---
        % Base del pivote
        \filldraw[fill=gray!30, draw=black, thick] (0,0) circle (0.22);
        % Anillo interior
        \filldraw[fill=white, draw=black, thin] (0,0) circle (0.12);
        % Pasador central
        \filldraw[fill=black] (0,0) circle (0.05);

        % --- 5. Ángulo y Velocidad Angular ---
        \draw[->, thick] (2.2, 0) arc (0:\Ang:2.2) node[midway, right, xshift=2pt] {$\omega_0$};

        % --- 6. Vector Gravedad ---
        \draw[->, thick] (4.5, 3.2) -- (4.5, 2.2) node[midway, right] {$\mathbf{g}$};

    \end{tikzpicture}
    \caption{Esquema del movimiento relativo de una partícula en el interior de un tubo rotatorio.}
    \label{fig:particula_tubo_giratorio}
\end{figure}